\section{Successive cuts and reoptimization}\label{cv:section}

A corner bound measures improvement in the current cone. After a cut is
added, reoptimization can select a different vertex and a different cone.
Thus local bound attainment does not determine the progress of a sequence of
LP relaxations. We first give a geometric condition that does guarantee
convergence. We then show why a ratio of local objective bounds cannot
replace that condition, and construct a sequence of maximal sets whose LP
values remain strictly below the nonlinear optimum.

\subsection{The loop and a sufficient depth condition}

Let $P_0=\{v\in\R^n:Av\le b\}$ be a polytope contained in
$[-B,B]^n$, and let $E$ be a finite set of bilinear terms. For
$e=(i,j,k)\in E$, put
\[
 S_e=\{v:v_k=v_iv_j\},\qquad
 F=P_0\cap\bigcap_{e\in E}S_e\ne\varnothing,
 \qquad z^*=\min_{v\in F}c^Tv.
\]
The value $z^*$ also equals the minimum over $\conv F$. In round $r$, solve
the LP over $P_r$ and choose an optimal vertex $v^r$ with an optimal
nonsingular basis of tight rows. Its basis cone is
\[
 K_r=v^r+\cone\{r_1^r,\ldots,r_n^r\},\qquad P_r\subseteq K_r.
\]
If $v^r\in F$, stop. Otherwise, for every violated term choose a closed
convex set containing $v^r$ in its interior and free of the side of the
bilinear inequality that $v^r$ violates. Add its intersection-cut halfspace
$H_{r,e}$ and retain every earlier cut. The sets may be defined in the
three coordinates of the term and lifted to $\R^n$. Every added cut is
valid for $P_r\cap S_e$, so
\[
 F\subseteq P_{r+1}\subseteq P_r\subseteq P_0.
\]
In particular, $c^Tv^r$ is nondecreasing and bounded above by $z^*$.
This is an exact loop: it does not discard cuts or stop at a positive
feasibility tolerance.

For a cut with steps $\alpha_j\in(0,\infty]$, define
\begin{equation}\label{cv:depth-definitions}
 d_{r,e}=\dist(v^r,K_r\cap H_{r,e}),\qquad
 \delta_{r,e}=\min_{j:\alpha_j<\infty}\alpha_j\norm{r_j^r},\qquad
 \gamma_r=\dist\left(0,\conv\left\{
                  \frac{r_j^r}{\norm{r_j^r}}:1\le j\le n\right\}\right).
\end{equation}
The quantity $d_{r,e}$ is the distance to the retained part of the cone,
which can exceed the distance to the cut hyperplane. The basis cone is
pointed, so $\gamma_r>0$. A positive lower bound independent of $r$ is a
stronger condition. At least one step is finite: if every step were
infinite, convexity and the interior vertex would imply
$K_r\subseteq\intt C$, contradicting the nonempty feasible set in $K_r$.

\begin{theorem}[Convergence under uniform cut depth]\label{cv:uniform-depth}
Suppose that for every $t>0$ there is $\eta(t)>0$ such that, whenever
$\dist(v^r,S_e)\ge t$, the cut added for term $e$ has
$d_{r,e}\ge\eta(t)$. Then the loop either stops at an optimal point of $F$,
or its LP values converge to $z^*$ and every accumulation point of its
vertices belongs to $F$ and is optimal.
\end{theorem}

\begin{proof}
If the loop stops, its point is feasible and its LP value is at most $z^*$,
so it is optimal. Otherwise compactness of $P_0$ supplies an accumulation
point $\widehat v$ and a subsequence $v^{r_\ell}\to\widehat v$. If
$\widehat v\notin S_e$, closedness of $S_e$ gives
$t=\dist(\widehat v,S_e)/2>0$. For all sufficiently large $\ell$,
$\dist(v^{r_\ell},S_e)\ge t$, and round $r_\ell$ adds a cut of depth at
least $\eta(t)$. For $m>\ell$, retention of that cut gives
\[
 v^{r_m}\in P_{r_m}\subseteq P_{r_\ell+1}
       \subseteq K_{r_\ell}\cap H_{r_\ell,e}.
\]
Consequently $\norm{v^{r_m}-v^{r_\ell}}\ge\eta(t)$, contradicting the
convergence of the subsequence. Thus every accumulation point is in $F$.
The monotone LP values have a limit $z\le z^*$, and continuity gives
$c^T\widehat v=z$. Feasibility implies $z\ge z^*$.
\end{proof}

The proof uses compactness, closed feasibility sets, valid retained cuts,
and a deep cut for each term that remains violated. It does not require
positive reduced costs, maximal sets, or bilinear structure. Those features
enter the following sufficient conditions for depth.

\begin{lemma}[Step length and cone geometry]\label{cv:cone-depth}
For every cut in \cref{cv:depth-definitions},
$d_{r,e}\ge\gamma_r\delta_{r,e}$.
\end{lemma}

\begin{proof}
Write $u_j=r_j^r/\norm{r_j^r}$ and
$\mu_j=\lambda_j\norm{r_j^r}$ for a point
$v^r+\sum_j\lambda_jr_j^r\in K_r\cap H_{r,e}$. The cut implies
\[
 1\le\sum_{j:\alpha_j<\infty}
             \frac{\mu_j}{\alpha_j\norm{r_j^r}}
       \le\frac{\sum_j\mu_j}{\delta_{r,e}}.
\]
Since $\sum_j\mu_j>0$, the definition of $\gamma_r$ also gives
\[
 \norm{\sum_j\mu_ju_j}
 \ge\gamma_r\sum_j\mu_j\ge\gamma_r\delta_{r,e}.
\]
Taking the infimum proves the claim.
\end{proof}

\begin{corollary}[Uniform pointedness and steps]\label{cv:pointed-convergence}
If $\gamma_r\ge\gamma_0>0$ in every round and, for every $t>0$,
$\delta_{r,e}\ge\eta_0(t)>0$ whenever $\dist(v^r,S_e)\ge t$, then the
conclusion of \cref{cv:uniform-depth} holds.
\end{corollary}

\begin{proof}
Apply \cref{cv:cone-depth} with $\eta(t)=\gamma_0\eta_0(t)$.
\end{proof}

\begin{corollary}[One selected term per round]\label{cv:one-term}
Instead add a cut for just one violated term $e_r$ per round. The conclusion
of \cref{cv:uniform-depth} still holds under either of these conditions:
\begin{enumerate}[label=(\alph*)]
\item $e_r$ maximizes $\dist(v^r,S_e)$, and the depth hypothesis holds for
      the cut actually added.
\item Each term has pairwise distinct indices, $e_r$ maximizes
      $q_e(v^r)=\abs{v_k^r-v_i^rv_j^r}$, $\gamma_r\ge\gamma_0>0$, and
      $\delta_{r,e_r}\ge\beta q_{e_r}(v^r)$ for a fixed $\beta>0$.
\end{enumerate}
\end{corollary}

\begin{proof}
Use the subsequence and the term $e$ in the proof of
\cref{cv:uniform-depth}. In (a), the selected distance is at least $t$,
so the cut added has depth at least $\eta(t)$. In (b), changing only the
coordinate $v_k^r$ to $v_i^rv_j^r$ produces a point of $S_e$, whence
$q_e(v^r)\ge\dist(v^r,S_e)$. The selected violation is therefore at least
$t$, and \cref{cv:cone-depth} gives depth at least $\gamma_0\beta t$.
Every later vertex satisfies the selected cut, so the same contradiction
applies.
\end{proof}

An arbitrary choice of a violated term does not ensure the cut needed in
this proof. The corollary makes no claim for a selection rule that can
indefinitely postpone a term.

\subsection{Rules that provide a step bound}

For this subsection the indices of each term are pairwise distinct. Write
$s=(x,y,w)$ for its three coordinates, let $\sbar$ be the projected LP
vertex, and put $q=\abs{\bar w-\bar x\bar y}>0$. A ball contained in the
chosen set gives a step bound: if $C\supseteq B(\sbar,\rho)$ and $p_j$ is
the projected ray, then
\[
 \alpha_j\norm{r_j^r}\ge\rho,
\]
because $\norm{p_j}\le\norm{r_j^r}$; a zero projected ray has infinite
step. A lower bound on steps still needs cone geometry to yield a lower
bound on cut depth.

\begin{lemma}[Uniform steps for bilinear rules]\label{cv:rule-steps}
Suppose $P_0\subseteq[-B,B]^n$, and set
\[
 L_B=\left(B^2+\frac{(B+1)^2}{4}\right)^{1/2},\qquad
 T_B=\sqrt{3B^2+1}.
\]
The following bounds hold for either violated side of a bilinear term:
\begin{enumerate}[label=(\alph*)]
\item The default Case-4 bilinear set of \citet{ChmielaMunozSerrano2020ZIB}, defined explicitly in
      \cref{cv:scip-definition}, gives
      $\delta_{r,e}\ge q/(2\sqrt2 L_B)$.
\item The orbit set with $F^T=M_\sigma(\sbar)^{-1}$ gives
      $\delta_{r,e}\ge q/T_B$, where
      \[
      M_+(s)=\begin{pmatrix}w&x\\y&1\end{pmatrix},\qquad
      M_-(s)=\begin{pmatrix}x&w\\1&y\end{pmatrix},
      \]
      and $\sigma$ is chosen so $\det M_\sigma(\sbar)=q>0$.
\item Fix $0<\theta\le1$. Suppose an idealized rule returns an orbit set
      whose value $\min_j u_j\alpha_j$ is at least $\theta$ times the
      supremum of that value over all orbit sets containing $\sbar$ in
      their interior. With $u_j=\norm{r_j^r}$, it gives
      $\delta_{r,e}\ge\theta q/T_B$. With
      $u_j=\norm{r_j^r}(\widehat\omega_j+\tau)$,
      $0\le\widehat\omega_j\le1$, and fixed $\tau>0$, it gives
      \[
      \delta_{r,e}\ge
             \frac{\theta\tau q}{(1+\tau)T_B}.
      \]
\end{enumerate}
The constants $\theta$ and $\tau$ are independent of the round and term.
\end{lemma}

The proof in \cref{cv:ball-proofs} exhibits a ball in each reference set
and compares the idealized objectives with the ball supplied by (b).
The two objectives in (c) correspond to the geometric-step rule and the
Euclidean reduced-cost rule with a positive perturbation, respectively.
Since $q\ge\dist(v^r,S_e)$, these bounds establish the step hypothesis of
\cref{cv:pointed-convergence}, and the residual bound in
\cref{cv:one-term}(b). Thus these rules converge under uniform pointedness.
The assertion for the optimized rules assumes a uniform multiplicative
accuracy $\theta$. A fixed absolute bisection error does not imply that
accuracy when the supremum can be small; it supplies no such guarantee for
a numerical implementation by itself.

\subsection{Exact local values can accompany small depth}

\begin{proposition}[Bound-optimal sets can be shallow]\label{cv:shallow-optimal}
Let $S=\{(x,y,w):w\le xy\}$, $\sbar=(0,0,1)$, and take the ordered rays
$e_y,e_x,e_w$ with positive reduced costs $\rc=(1,\varepsilon,1)$,
$0<\varepsilon\le1$. Then $\zK(\rc)=2\sqrt\varepsilon$.
Every closed convex $S$-free set containing $\sbar$ in its interior
satisfies $\alpha_y\alpha_x\le4$. Consequently, a set satisfying
$z_C(\rc)\ge\theta\zK(\rc)$, $0<\theta\le1$, has
\[
 \delta\le\frac{2\sqrt\varepsilon}{\theta},\qquad
 d\le\frac{2\sqrt\varepsilon}{\theta}.
\]
Moreover, maximal sets attain $z_C(\rc)=\zK(\rc)$ while $d\to0$, even
though $\dist(\sbar,S)=1$ and the cone is fixed and uniformly pointed.
\end{proposition}

\begin{proof}
A point in the cone is $(\lambda_x,\lambda_y,1+\lambda_w)$.
Its membership in $S$ requires
$\lambda_x\lambda_y\ge1+\lambda_w$. Minimization sets
$\lambda_w=0$, and AM--GM gives
\[
 \zK(\rc)=\min_{\lambda_x\lambda_y\ge1}
                 (\lambda_y+\varepsilon\lambda_x)=2\sqrt\varepsilon.
\]
For finite positive $t_y\le\alpha_y$ and $t_x\le\alpha_x$, the midpoint
of $(0,t_y,1)$ and $(t_x,0,1)$ belongs to $C$. If $t_yt_x>4$, moving
slightly from that midpoint toward the interior point $\sbar$ produces a
point in $\intt C\cap S$, a contradiction. Thus $t_yt_x\le4$ for all
such choices. Since both steps are positive, neither can be infinite,
and $\alpha_y\alpha_x\le4$. The assumed bound implies
$\varepsilon\alpha_x\ge2\theta\sqrt\varepsilon$, and hence
$\alpha_y\le2\sqrt\varepsilon/\theta$. The ray intercept is in the
retained cone, so $d\le\alpha_y$ as well.

For attainment, take
\[
 C_\varepsilon=\left\{w\ge
              \frac{(\sqrt\varepsilon x+y/\sqrt\varepsilon)^2}{4}\right\}.
\]
These are orbit sets with $F^T=\diag(\sqrt\varepsilon,1/\sqrt\varepsilon)$,
and they are maximal by \cref{cv:stalling}. Their steps are
$(2\sqrt\varepsilon,2/\sqrt\varepsilon,\infty)$ in the stated order,
so the local ratio is exactly one. The resulting halfspace is
$\sqrt\varepsilon x+y/\sqrt\varepsilon\ge2$, whose closest point to
$\sbar$ lies in the cone. Therefore
\[
 d=\frac{2\sqrt\varepsilon}{\sqrt{1+\varepsilon^2}}\longrightarrow0.
\]
If $w\le0$, a point of $S$ is at distance at least one from $\sbar$.
If $w>0$, then $x^2+y^2\ge2xy\ge2w$, so its squared distance is at least
$2w+(w-1)^2=1+w^2$. The point $(0,0,0)\in S$ attains distance one.
\end{proof}

The cone in this example has $\gamma=1/\sqrt3$. It is also the basis
cone at the unique LP optimum of the fixed box relaxation
$P_0=[0,2]^2\times[1,2]$ for the objectives
$\varepsilon x+y+w$; the nonlinear feasible set is nonempty. Thus fixed
box bounds and uniform pointedness do not convert the raw local ratio
$z_C/\zK$ into an objective-independent depth guarantee. The objective
changes with $\varepsilon$, and its smallest positive reduced cost tends
to zero. This proposition does not prove that a bound-optimal rule lacks
a depth modulus along one fixed problem, or that its loop fails to converge.

\subsection{Bound attainment and ties}

\begin{proposition}[Ties forced by a supported corner minimizer]\label{cv:dual-ties}
Let $\rc\in\Rpp^N$, let $S$ be closed, and suppose the corner minimum is
attained at $\lambda^*\ge0$, with
$\sbar+P\lambda^*\in S$ and
$\rc^T\lambda^*=\zK(\rc)\in(0,\infty)$. Put
$J=\supp\lambda^*$. If a valid intersection cut attains
$z_C(\rc)=\zK(\rc)$, then
\[
 \omega_j\alpha_j=\zK(\rc)\quad(j\in J),\qquad
 \lambda^*\in\relint\conv\{\alpha_je_j:j\in J\}.
\]
This simplex belongs to the optimal face of the corner LP after the cut.
If $|J|\ge2$, each intercept basis in that simplex has an additional
zero reduced cost. Thus the corner LP has a dual tie and its optimal face
contains a feasible corner point.
\end{proposition}

\begin{proof}
Validity and bound attainment give
\[
 1\le\sum_{j\in J}\frac{\lambda_j^*}{\alpha_j}
    \le\sum_{j\in J}\frac{\omega_j\lambda_j^*}{\zK(\rc)}=1.
\]
Every summand with $j\in J$ has positive coefficient, so equality forces
$1/\alpha_j=\omega_j/\zK(\rc)$. The positive numbers
$\lambda_j^*/\alpha_j$ sum to one and express $\lambda^*$ as a relative
interior convex combination of the intercepts. Every intercept has value
$\zK(\rc)$, the minimum of the corner LP. At an intercept basis indexed
by $i\in J$, the reduced cost of another variable $j\in J$ is
$\omega_j-(\zK(\rc)/\alpha_j)=0$.
\end{proof}

Only ties on the support of $\lambda^*$ are forced. The proposition does
not imply ties when that support has size one, and it concerns the corner
LP, whose intercepts need not belong to the full current relaxation. It
does not establish a cause of slower reoptimization.

\subsection{A suboptimal limit with maximal sets}

\begin{theorem}[Stalling on uniformly pointed cones]\label{cv:stalling}
Fix $0<\varepsilon<1$, $X\ge2/\sqrt\varepsilon$,
$Y\ge\sqrt\varepsilon$, and $W>1$. For
\[
 \min\{\varepsilon x+y+w:
     0\le x\le X,\ 0\le y\le Y,\ 1\le w\le W,\ w=xy\},
\]
the optimum is $z^*=1+2\sqrt\varepsilon$. For $a>0$, every set
\[
 C_a=\{(x,y,w):w\ge(ax+y/a)^2/4\}
\]
is maximal for $S=\{w\le xy\}$ and is the image of the default Case-4
set at $(0,0,1)$ under $(x,y,w)\mapsto(x/a,ay,w)$.
Given any strictly increasing sequence
\[
 0=\xi_0<\xi_1<\cdots\longrightarrow\xi_\infty<2/\sqrt\varepsilon,
\]
start from the box relaxation and use $C_{2/\xi_{r+1}}$ in round $r$.
The unique LP optimum in round $r$ is
\[
 v^r=(\xi_r,0,1),\qquad c^Tv^r=1+\varepsilon\xi_r.
\]
The cones have a uniform positive pointedness bound, but the cut depths
tend to zero. The LP values converge to
$1+\varepsilon\xi_\infty<z^*$, although every vertex has violation one.
The same conclusion holds when the nonlinear constraint is $w\le xy$.
\end{theorem}

The construction keeps its steps along one ray arbitrarily small. Its
cuts have the explicit form
\[
 \frac{x}{\xi_{r+1}}+\frac{\xi_{r+1}y}{4}\ge1.
\]
The next vertex moves only from $(\xi_r,0,1)$ to $(\xi_{r+1},0,1)$.
\Cref{cv:stalling-proof} proves maximality, all basis and cut identities,
and a uniform separating-vector bound for the cones. For example,
$\varepsilon=1/4$ and $\xi_r=2(1-2^{-r})$ give limit $3/2$ and optimum
$2$. Any value strictly between the initial LP value and $z^*$ can be
obtained by choosing $\xi_\infty$.

At the initial corner, in contrast, the bound-optimal set
$C_{\sqrt\varepsilon}$ gives the cut
$\varepsilon x+y\ge2\sqrt\varepsilon$. Together with $w\ge1$ this
already gives the true optimum, which the feasible point
$(1/\sqrt\varepsilon,\sqrt\varepsilon,1)$ attains. Thus the example
also admits a bound-optimal choice that closes the gap in one round.

This is an adversarial choice among maximal sets. The sets are images of
the initial default set; the construction does not follow the default
rule at subsequent vertices. Nor does it require the initial relaxation
to contain McCormick inequalities. It proves failure for an unrestricted
maximal-set selection rule in the stated loop, while convergence of the
natural default and bound-optimal rules without the sufficient geometric
hypotheses remains open.
